\documentclass[11pt]{article}
\usepackage[a4paper,margin=1in]{geometry}
\usepackage{fontspec}
\usepackage[boldfont=false]{xeCJK}
\setCJKmainfont{FandolSong-Regular.otf}
\usepackage{amsmath}
\usepackage{amssymb}
\usepackage{array}
\usepackage{booktabs}
\usepackage{graphicx}
\usepackage{adjustbox}
\usepackage{enumitem}
\setlist[itemize]{topsep=2pt,partopsep=0pt,parsep=0pt,itemsep=2pt,leftmargin=1.4em}
\setlist[enumerate]{topsep=2pt,partopsep=0pt,parsep=0pt,itemsep=2pt,leftmargin=1.6em}
\usepackage{parskip}
\usepackage{fvextra}
\fvset{breaklines=true,breakanywhere=true,fontsize=\small}
\setlength{\parindent}{0pt}

\begin{document}

\begin{center}
  {\LARGE\bfseries Economic development}\\[0.35em]
  {\large Igcse Economics}
\end{center}
\vspace{0.6em}

\section*{Living standards}

"Living standards" means how well people live — their income, health, education, and quality of life.

\subsection*{Indicators of living standards}

An \textbf{indicator} 指标 is a number we use to measure something. Two indicators of \textbf{living standards} 生活水平 are common.

\textbf{Real GDP per head} is the main one. \textbf{GDP} (\textbf{gross domestic product} 国内生产总值) is the total value of goods and services a country makes in a year. \textbf{Real} 实际 means after rising prices are taken out. \textbf{Per head} 人均 means divided by the number of people. So real GDP per head shows the average output — and roughly the average \textbf{income} 收入 — of each person. Higher real GDP per head usually means higher living standards.

The \textbf{Human Development Index} 人类发展指数 (HDI) is a wider measure. It combines three things: income per head, education (years of schooling), and health (measured by \textbf{life expectancy} 预期寿命 — how long people are expected to live). HDI gives a fuller picture than GDP alone.

\subsection*{Limits of the comparison}

Comparing living standards is not simple:

\begin{itemize}
\item GDP per head is an \emph{average}. It hides the \textbf{income distribution} 收入分配 — a country can have a high average but a few rich people and many poor people.

\item it leaves out unpaid work, and work in the hidden economy.

\item it does not measure pollution, free time, or safety.

\item prices and money differ between countries, so the same income buys different amounts.

\end{itemize}

This is why HDI, which adds health and education, is often a better guide.

\section*{Poverty}

\textbf{Poverty} 贫困 is when people cannot afford a basic standard of living. There are two types:

\begin{itemize}
\item \textbf{absolute poverty} 绝对贫困 — people cannot afford even the basics: enough food, clean water, shelter, and clothing.

\item \textbf{relative poverty} 相对贫困 — people are poor \emph{compared with others} in the same country. They can afford the basics but much less than most people around them.

\end{itemize}

\subsection*{Causes of poverty}

\begin{itemize}
\item low income, low \textbf{wages} 工资, or no job (unemployment).

\item poor education and poor health, which make it hard to get good work.

\item old age, or having many children to support.

\end{itemize}

\subsection*{Policies to reduce poverty}

\begin{itemize}
\item \textbf{economic growth} 经济增长 — a growing economy creates jobs and income.

\item \textbf{redistribute} 再分配 income — use progressive taxes on the rich to pay for help for the poor.

\item a \textbf{minimum wage} 最低工资 — a legal lowest wage, so low-paid workers earn more.

\item state \textbf{benefits} 福利 — money paid by the government to the poor, old, or unemployed.

\item better state education and healthcare, so the poor can improve their lives.

\end{itemize}

\section*{Population}

A country's \textbf{population} 人口 is the number of people living in it. Three things change its size.

\begin{itemize}
\item the \textbf{birth rate} 出生率 — the number of births each year for every 1000 people.

\item the \textbf{death rate} 死亡率 — the number of deaths each year for every 1000 people.

\item \textbf{migration} 移民 — people moving between countries. \textbf{immigration} 迁入 is people moving \emph{in}; \textbf{emigration} 迁出 is people moving \emph{out}. \textbf{Net migration} 净移民 is immigration minus emigration.

\end{itemize}

The population grows when births plus immigration are more than deaths plus emigration.

\subsection*{Population structure}

The \textbf{population structure} 人口结构 is how the population is split by age and sex. We often draw it as a \textbf{population pyramid} 人口金字塔 — a chart with age groups stacked up, males on one side and females on the other.

\subsection*{An ageing population}

In many richer countries the \textbf{ageing population} 人口老龄化 is a big change: people live longer and have fewer children, so the share of old people rises. Effects:

\begin{itemize}
\item fewer workers compared with the number of retired people, so each worker must support more people.

\item the government spends more on pensions and healthcare.

\item it may need higher taxes, or workers from other countries, to fill the gap.

\end{itemize}

\section*{Differences in economic development between countries}

\textbf{Economic development} 经济发展 means a rise in people's well-being, not just more output. We often split the world into \textbf{developed countries} 发达国家 (rich, like Japan or Germany) and \textbf{developing countries} 发展中国家 (poorer, with lower incomes).

\subsection*{Why countries differ}

\begin{itemize}
\item \textbf{income} — developed countries have much higher income per head.

\item \textbf{productivity} 生产率 — developed countries use more capital and skill, so each worker produces more. In many developing countries, farmers do \textbf{subsistence farming} 自给农业 (growing food just for their own family), which produces little to sell.

\item \textbf{population} — developing countries often have higher birth rates and younger populations.

\item \textbf{resources, education, and health} — better schools, hospitals, and use of resources all raise development.

\item \textbf{debt} 债务 — many developing countries owe large amounts of money, and paying it back leaves less for schools and hospitals.

\item \textbf{trade} 贸易 — developed countries trade more, and often sell higher-value goods.

\end{itemize}

\subsection*{How developed and developing countries are linked}

The two groups depend on each other. Developing countries sell raw materials and cheap goods to developed countries, and buy machines and skills from them. Trade, aid, and investment from richer countries can help poorer countries develop — but developing countries can also become too dependent on them.


\end{document}
